% !TeX root = ../教材题库_全部条目_转写工作稿.tex
% Source: XJ-B1-C03-EX-P066-Y0934
% 第3章 / 3.1 函数 / 教材 p.66 / PDF 第73页
% OCR未校对：公式与文字均可能有误；未完成数学TeX转写；禁止直接发给学生
% 原题及相邻原解截图（核对用）：../题目原图/XJ-B1-C03-EX-P066-Y0934-01.webp
% 本文件只在转写完成后作为题干使用；OCR 条目尚不是可靠题干。
\begin{bankitem}{3.1 例题 例题1}{第3章 / 3.1 函数 / 教材 p.66 / PDF 第73页}{XJ-B1-C03-EX-P066-Y0934}{OCR未校对：公式与文字均可能有误；未完成数学TeX转写；禁止直接发给学生}
\begin{ocrdraft}
[教材p.66 / PDF第73页]
\par
00确定下列函数的定义域：
\par
（D不）一十x一2；
\par
工2十3工十2．
\par
解（1）要使二次根式有意义，实数工需满足主2十x一2冫0，即
\par
一2或冫
\par
1.故函数不到的定义域为（一。。，一2]UCI，十。。）．
\par
（2）要使分式有意义，当且仅当分式的分母十十2一
\par
(x十1）0十2）乒0，即需满足x乒一1且工乒一2．故函数g(x）
\par
的定义域为（一。。
\par
一2）U（一2，
\par
十。。）；也可
\par
即使两个函数的对应
\par
表示为00学一1且一2\}或R\textbackslash{}\{
\par
关系在形式上相司，但定
\par
义域不司，它们也不是司
\par
两个函数不和）是同一个函数，当且仅当它们有相
\par
一个函数．
\par
同的定义域U且对每个\ensuremath{\text{[OCR:EUR]}}LJ都有丆0）一g(x）．
\end{ocrdraft}
\end{bankitem}
